% Part: second-order-logic
% Chapter: syntax-and-semantics
% Section: satisfaction

\documentclass[../../../include/open-logic-section]{subfiles}

\begin{document}

\olfileid{sol}{syn}{sat}

\olsection{الإرضاء}

\begin{explain}
تعريف علاقة الإرضاء~$\Sat{\OLId{latin-looped}{M}}{!A}[\OLId{latin-ordinary}{s}]$ لـ!!{formula}s الرتبة الثانية
يقتضي أن تشمل تعريفاتنا !!{variable}s هذه الرتبة. أما مفهوم
!!a{structure} فلا يتغير بالانتقال من الرتبة الأولى إلى الثانية.
وإنما يتغير تعيين المتغيرات~$\OLId{latin-ordinary}{s}$؛ فلا يكفي فيه إسناد القيم إلى
!!{variable}s الرتبة الأولى، بل لا بد من إسناد ما يلائم
!!{variable}s الرتبة الثانية أيضًا.
\end{explain}

\begin{defn}[تعيين المتغيرات]
\emph{تعيين المتغيرات}~$\OLId{latin-ordinary}{s}$ لـ!!a{structure}~$\Struct{M}$ دالة
تختص كل متغير بما يأتي:
\begin{enumerate}
\item !!{variable} الأفراد~$\Obj{v_i}$ بعنصر من~$\Domain \OLId{latin-looped}{M}$، فتكون
  $\OLId{latin-ordinary}{s}(\Obj{v_i}) \in \Domain{\OLId{latin-looped}{M}}$؛
\item متغير العلاقة ذي~$\OLId{latin-ordinary}{n}$ مواضع~$\Obj{V_i^n}$ بعلاقة ذات~$\OLId{latin-ordinary}{n}$
  مواضع على~$\Domain{\OLId{latin-looped}{M}}$، فتكون
  $\OLId{latin-ordinary}{s}(\Obj{V_i^n}) \subseteq \Domain{\OLId{latin-looped}{M}}^\OLId{latin-ordinary}{n}$؛
\item متغير الدالة ذي~$\OLId{latin-ordinary}{n}$ مواضع~$\Obj{u_i^n}$ بدالة ذات~$\OLId{latin-ordinary}{n}$ مواضع
  من~$\Domain{\OLId{latin-looped}{M}}^\OLId{latin-ordinary}{n}$ إلى~$\Domain{\OLId{latin-looped}{M}}$، أي
  $\OLId{latin-ordinary}{s}(\Obj{u_i^n})\colon \Domain{\OLId{latin-looped}{M}}^\OLId{latin-ordinary}{n} \to \Domain{\OLId{latin-looped}{M}}$؛
\end{enumerate}
\end{defn}

\begin{explain}
تعطي~!!^a{structure} !!a{value} لكل~!!{constant} و!!{function}،
وعلاقة لكل محمول. وأما تعيين متغيرات الرتبة الثانية فيعطي متغيرات
الأفراد أفرادًا، ومتغيرات الدوال دوال، ومتغيرات العلاقات علاقات.
فإذا اجتمع هذان أمكن تعيين قيمة كل حد.
\end{explain}

\begin{defn}[\usetoken{S}{value} حد]
ليكن~$\OLId{latin-ordinary}{t}$ حدًا من اللغة~$\Lang L$، ولتكن~$\Struct M$
!!a{structure} للغة~$\Lang L$، وليكن~$\OLId{latin-ordinary}{s}$ تعيين~!!a{variable}
لـ$\Struct M$. نعرّف \emph{!!{value}}~$\Value{\OLId{latin-ordinary}{t}}{\OLId{latin-looped}{M}}[\OLId{latin-ordinary}{s}]$ ببنود
تعريف قيمة حدود الرتبة الأولى، ونزيد عليها هذا البند:
\begin{quote}
\indcase{t}{\Atom{u}{t_1, \ldots, t_n}}{
\[
\Value{\indfrm}{\OLId{latin-looped}{M}}[\OLId{latin-ordinary}{s}] = \OLId{latin-ordinary}{s}(\OLId{latin-ordinary}{u})(\Value{\OLId{latin-ordinary}{t}_\OLMathNumeral{1}}{\OLId{latin-looped}{M}}[\OLId{latin-ordinary}{s}], \ldots,
\Value{\OLId{latin-ordinary}{t}_\OLId{latin-ordinary}{n}}{\OLId{latin-looped}{M}}[\OLId{latin-ordinary}{s}]).
\]}
\end{quote}
\end{defn}

\begin{defn}[المغاير في~$\OLId{latin-ordinary}{x}$]
ليكن~$\OLId{latin-ordinary}{s}$ تعيين~!!a{variable} لـ!!a{structure}~$\Struct M$.
كل تعيين~!!{variable}~$\OLId{latin-ordinary}{s}'$ لـ$\Struct M$ يوافق~$\OLId{latin-ordinary}{s}$ في جميع
المواضع، ولا يجوز أن يخالفه إلا في القيمة المسندة إلى~$\OLId{latin-ordinary}{x}$، يسمى
\emph{مغايرًا في~$\OLId{latin-ordinary}{x}$} لـ$\OLId{latin-ordinary}{s}$. ومتى كان
$\OLId{latin-ordinary}{s}'$ مغايرًا في~$\OLId{latin-ordinary}{x}$ لـ$\OLId{latin-ordinary}{s}$ رمزنا إلى ذلك
بـ$\varAssign{\OLId{latin-ordinary}{s}'}{\OLId{latin-ordinary}{s}}{\OLId{latin-ordinary}{x}}$. ولا نشترط وقوع الاختلاف فعلًا.
وكذلك يعرّف المغاير في متغير الرتبة الثانية~$\OLId{latin-looped}{X}$ أو~$\OLId{latin-ordinary}{u}$.
\end{defn}

\begin{defn}
  ليكن~$\OLId{latin-ordinary}{s}$ تعيين~!!a{variable} لـ!!a{structure}~$\Struct M$،
  وليكن~$\OLId{latin-ordinary}{m} \in \Domain{\OLId{latin-looped}{M}}$. نرمز بـ$\Subst{\OLId{latin-ordinary}{s}}{\OLId{latin-ordinary}{m}}{\OLId{latin-ordinary}{x}}$ إلى
  تعيين~!!{variable} الذي يرد تعريفه في المعادلة:
  \[\Subst{\OLId{latin-ordinary}{s}}{\OLId{latin-ordinary}{m}}{\OLId{latin-ordinary}{x}}(\OLId{latin-ordinary}{y}) = \begin{cases}
    \OLId{latin-ordinary}{m} & \text{إذا كان } \OLId{latin-ordinary}{y} \ident \OLId{latin-ordinary}{x}\\
    \OLId{latin-ordinary}{s}(\OLId{latin-ordinary}{y}) & \text{فيما عدا ذلك،}
  \end{cases}\]
  وليكن~$\OLId{latin-looped}{X}$ !!{variable} علاقة ذا~$\OLId{latin-ordinary}{n}$ مواضع، ولتكن~$\OLId{latin-looped}{R} \subseteq
  \Domain{\OLId{latin-looped}{M}}^\OLId{latin-ordinary}{n}$. عندئذ يكون~$\Subst{\OLId{latin-ordinary}{s}}{\OLId{latin-looped}{R}}{\OLId{latin-looped}{X}}$ تعيين~!!{variable}
  المحدد بالمعادلة:
  \[\Subst{\OLId{latin-ordinary}{s}}{\OLId{latin-looped}{R}}{\OLId{latin-looped}{X}}(\OLId{latin-ordinary}{y}) = \begin{cases}
    \OLId{latin-looped}{R} & \text{إذا كان } \OLId{latin-ordinary}{y} \ident \OLId{latin-looped}{X}\\
    \OLId{latin-ordinary}{s}(\OLId{latin-ordinary}{y}) & \text{فيما عدا ذلك}.
  \end{cases}\]
  وأما إذا كان~$\OLId{latin-ordinary}{u}$ !!{variable} دالة ذا~$\OLId{latin-ordinary}{n}$ مواضع، وكانت
  $\OLId{latin-ordinary}{f}\colon \Domain{\OLId{latin-looped}{M}}^\OLId{latin-ordinary}{n} \to \Domain{\OLId{latin-looped}{M}}$، فإن~$\Subst{\OLId{latin-ordinary}{s}}{\OLId{latin-ordinary}{f}}{\OLId{latin-ordinary}{u}}$
  تعيين~!!{variable} المحدد بالمعادلة:
  \[\Subst{\OLId{latin-ordinary}{s}}{\OLId{latin-ordinary}{f}}{\OLId{latin-ordinary}{u}}(\OLId{latin-ordinary}{y}) = \begin{cases}
    \OLId{latin-ordinary}{f} & \text{إذا كان } \OLId{latin-ordinary}{y} \ident \OLId{latin-ordinary}{u}\\
    \OLId{latin-ordinary}{s}(\OLId{latin-ordinary}{y}) & \text{فيما عدا ذلك}.
  \end{cases}\]
  ولا يختص~$\OLId{latin-ordinary}{y}$ في هذه الحالات برتبة؛ فيجوز أن يكون أي~!!{variable}
  من الرتبة الأولى أو الثانية.
\end{defn}

\begin{defn}[الإرضاء]
نعرف إرضاء !!{formula}s الرتبة الثانية~$!A$ ببنود
\olref[fol][syn][sat]{defn:satisfaction}، ونضيف إليها:
\begin{enumerate}
\item \indcase{!A}{\Atom{X^n}{t_1, \dots, t_n}}{$\Sat{\OLId{latin-looped}{M}}{\indfrm}[\OLId{latin-ordinary}{s}]$
  إذا وفقط إذا كان
  $\langle \Value{\OLId{latin-ordinary}{t}_\OLMathNumeral{1}}{\OLId{latin-looped}{M}}[\OLId{latin-ordinary}{s}], \dots, \Value{\OLId{latin-ordinary}{t}_\OLId{latin-ordinary}{n}}{\OLId{latin-looped}{M}}[\OLId{latin-ordinary}{s}] \rangle \in
  \OLId{latin-ordinary}{s}(\OLId{latin-looped}{X}^\OLId{latin-ordinary}{n})$.}
\tagitem{prvAll}{%
  \indcase{!A}{\lforall[X][!B]}{$\Sat{\OLId{latin-looped}{M}}{\indfrm}[\OLId{latin-ordinary}{s}]$ إذا وفقط إذا
  كان لكل~$\OLId{latin-looped}{R} \subseteq \Domain{\OLId{latin-looped}{M}}^\OLId{latin-ordinary}{n}$
  $\Sat{\OLId{latin-looped}{M}}{!B}[\Subst{\OLId{latin-ordinary}{s}}{\OLId{latin-looped}{R}}{\OLId{latin-looped}{X}}]$.}}{}

\tagitem{prvEx}{%
  \indcase{!A}{\lexists[X][!B]}{$\Sat{\OLId{latin-looped}{M}}{\indfrm}[\OLId{latin-ordinary}{s}]$ إذا وفقط إذا
  وجد~$\OLId{latin-looped}{R} \subseteq \Domain{\OLId{latin-looped}{M}}^\OLId{latin-ordinary}{n}$ واحد على الأقل يحقق
  $\Sat{\OLId{latin-looped}{M}}{!B}[\Subst{\OLId{latin-ordinary}{s}}{\OLId{latin-looped}{R}}{\OLId{latin-looped}{X}}]$.}}{}

\tagitem{prvAll}{%
  \indcase{!A}{\lforall[u][!B]}{$\Sat{\OLId{latin-looped}{M}}{\indfrm}[\OLId{latin-ordinary}{s}]$ إذا وفقط إذا
  كان لكل~$\OLId{latin-ordinary}{f}\colon \Domain{\OLId{latin-looped}{M}}^\OLId{latin-ordinary}{n} \to \Domain{\OLId{latin-looped}{M}}$
  $\Sat{\OLId{latin-looped}{M}}{!B}[\Subst{\OLId{latin-ordinary}{s}}{\OLId{latin-ordinary}{f}}{\OLId{latin-ordinary}{u}}]$.}}{}

\tagitem{prvEx}{%
  \indcase{!A}{\lexists[u][!B]}{$\Sat{\OLId{latin-looped}{M}}{\indfrm}[\OLId{latin-ordinary}{s}]$ إذا وفقط إذا
  وجدت~$\OLId{latin-ordinary}{f}\colon \Domain{\OLId{latin-looped}{M}}^\OLId{latin-ordinary}{n} \to \Domain{\OLId{latin-looped}{M}}$ واحدة على الأقل تحقق
  $\Sat{\OLId{latin-looped}{M}}{!B}[\Subst{\OLId{latin-ordinary}{s}}{\OLId{latin-ordinary}{f}}{\OLId{latin-ordinary}{u}}]$.}}{}
\end{enumerate}
\end{defn}

\begin{ex}
  لننظر في !!{formula}~$\lforall[\OLId{latin-ordinary}{z}][(\Atom{X}{z} \liff \lnot
    \Atom{Y}{z})]$. ليس فيها تكميم من الرتبة الثانية، وإن كانت
  !!{variable}s الرتبة الثانية~$\OLId{latin-looped}{X}$ و$\OLId{latin-looped}{Y}$ واردة فيها؛ ونعد كلا
  المتغيرين هنا أحادي الموضع. وأما !!{sentence} المناظرة من الرتبة
  الأولى
  $\lforall[\OLId{latin-ordinary}{z}][(\Atom{P}{z} \liff \lnot \Atom{R}{z})]$ فمؤداها أن
  ما يندرج تحت تأويل~$\OLId{latin-looped}{P}$ هو بعينه ما لا يندرج تحت تأويل~$\OLId{latin-looped}{R}$ في المجال.
  ففي~!!a{structure} يتحدد تأويل~!!a{predicate}~$\OLId{latin-looped}{P}$
  بالتأويل~$\Assign{\OLId{latin-looped}{P}}{\OLId{latin-looped}{M}}$. وأما !!{variable}s الرتبة الثانية
  من قبيل~$\OLId{latin-looped}{X}$ و$\OLId{latin-looped}{Y}$ فلا تتولى !!{structure} نفسها تأويلها، بل يتولاه
  تعيين~!!a{variable}. لذلك ليست !!{formula} الرتبة الثانية هذه
  !!a{sentence}؛ إذ تشتمل على !!{variable}s حرة هي~$\OLId{latin-looped}{X}$ و$\OLId{latin-looped}{Y}$.
  ولا يكفي في إرضائها ذكر !!a{structure}~$\Struct{M}$، حتى يضم
  إليها تعيين~!!a{variable}~$\OLId{latin-ordinary}{s}$.

  يكون~$\Sat{\OLId{latin-looped}{M}}{\lforall[\OLId{latin-ordinary}{z}][(\Atom{X}{z} \liff \lnot
    \Atom{Y}{z})]}[\OLId{latin-ordinary}{s}]$ حين تكون !!{element}s~$\OLId{latin-ordinary}{s}(\OLId{latin-looped}{X})$ هي بالضبط
  ما ليس من !!{element}s~$\OLId{latin-ordinary}{s}(\OLId{latin-looped}{Y})$، والعكس؛ أي إذا وفقط إذا
  $\OLId{latin-ordinary}{s}(\OLId{latin-looped}{Y}) = \Domain{\OLId{latin-looped}{M}} \setminus \OLId{latin-ordinary}{s}(\OLId{latin-looped}{X})$. ولنأخذ
  $\Domain{\OLId{latin-looped}{M}} = \{\OLMathNumeral{1}, \OLMathNumeral{2}, \OLMathNumeral{3}\}$. وليس في الصيغة !!{predicate}s
  ولا !!{function}s ولا !!{constant}s؛ فلا يعنينا من
  $\Struct{M}$ هنا إلا مجالها.
  فإذا كان~$\OLId{latin-ordinary}{s}_\OLMathNumeral{1}(\OLId{latin-looped}{X}) = \{\OLMathNumeral{1}, \OLMathNumeral{2}\}$ و$\OLId{latin-ordinary}{s}_\OLMathNumeral{1}(\OLId{latin-looped}{Y}) = \{\OLMathNumeral{3}\}$ حصل
  $\Sat{\OLId{latin-looped}{M}}{\lforall[\OLId{latin-ordinary}{z}][(\Atom{X}{z}
      \liff \lnot \Atom{Y}{z})]}[\OLId{latin-ordinary}{s}_\OLMathNumeral{1}]$.

  وأما إن كان~$\OLId{latin-ordinary}{s}_\OLMathNumeral{2}(\OLId{latin-looped}{X}) = \{\OLMathNumeral{1}, \OLMathNumeral{2}\}$ و
  $\OLId{latin-ordinary}{s}_\OLMathNumeral{2}(\OLId{latin-looped}{Y}) = \{\OLMathNumeral{2}, \OLMathNumeral{3}\}$، فالحال بخلاف ذلك:
  $\Sat/{\OLId{latin-looped}{M}}{\lforall[\OLId{latin-ordinary}{z}][(\Atom{X}{z} \liff \lnot
      \Atom{Y}{z})]}[\OLId{latin-ordinary}{s}_\OLMathNumeral{2}]$. ووجهه أن
  $\Sat{\OLId{latin-looped}{M}}{\Atom{X}{z}}[\Subst{\OLId{latin-ordinary}{s}_\OLMathNumeral{2}}{\OLMathNumeral{2}}{\OLId{latin-ordinary}{z}}]$، لأن
  $\OLMathNumeral{2} \in \Subst{\OLId{latin-ordinary}{s}_\OLMathNumeral{2}}{\OLMathNumeral{2}}{\OLId{latin-ordinary}{z}}(\OLId{latin-looped}{X})$، في حين أن
  $\Sat/{\OLId{latin-looped}{M}}{\lnot \Atom{Y}{z}}[\Subst{\OLId{latin-ordinary}{s}_\OLMathNumeral{2}}{\OLMathNumeral{2}}{\OLId{latin-ordinary}{z}}]$، لأن
  $\OLMathNumeral{2} \in \Subst{\OLId{latin-ordinary}{s}_\OLMathNumeral{2}}{\OLMathNumeral{2}}{\OLId{latin-ordinary}{z}}(\OLId{latin-looped}{Y})$ أيضًا.
\end{ex}

\begin{ex}
  يكون~$\Sat{\OLId{latin-looped}{M}}{\lexists[\OLId{latin-looped}{Y}][(\lexists[\OLId{latin-ordinary}{y}][\Atom{Y}{y}] \land
  \lforall[\OLId{latin-ordinary}{z}][(\Atom{X}{z} \liff \lnot \Atom{Y}{z})])]}[\OLId{latin-ordinary}{s}]$ إذا
  وجد~$\OLId{latin-looped}{N} \subseteq \Domain{\OLId{latin-looped}{M}}$ يحقق
  $\Sat{\OLId{latin-looped}{M}}{(\lexists[\OLId{latin-ordinary}{y}][\Atom{Y}{y}] \land \lforall[\OLId{latin-ordinary}{z}][(\Atom{X}{z}
  \liff \lnot \Atom{Y}{z})])}[\Subst{\OLId{latin-ordinary}{s}}{\OLId{latin-looped}{N}}{\OLId{latin-looped}{Y}}]$.
  ويتحقق ذلك متى اجتمع شرطان: أن يكون
  $\OLId{latin-looped}{N} \neq \emptyset$، ليحصل
  $\Sat{\OLId{latin-looped}{M}}{\lexists[\OLId{latin-ordinary}{y}][\Atom{Y}{y}]}[\Subst{\OLId{latin-ordinary}{s}}{\OLId{latin-looped}{N}}{\OLId{latin-looped}{Y}}]$،
  وأن يكون~$\OLId{latin-looped}{N} = \Domain{\OLId{latin-looped}{M}} \setminus \OLId{latin-ordinary}{s}(\OLId{latin-looped}{X})$ كما سبق.
  فحاصل القول أن
  $\Sat{\OLId{latin-looped}{M}}{\lexists[\OLId{latin-looped}{Y}][(\lexists[\OLId{latin-ordinary}{y}][\Atom{Y}{y}] \land
  \lforall[\OLId{latin-ordinary}{z}][(\Atom{X}{z} \liff \lnot \Atom{Y}{z})])]}[\OLId{latin-ordinary}{s}]$ إذا
  وفقط إذا كانت~$\Domain{\OLId{latin-looped}{M}} \setminus \OLId{latin-ordinary}{s}(\OLId{latin-looped}{X})$ غير خالية، أي
  $\OLId{latin-ordinary}{s}(\OLId{latin-looped}{X}) \neq \Domain{\OLId{latin-looped}{M}}$. ومن أمثلة إرضاء !!{formula} أن يكون
  $\Domain{\OLId{latin-looped}{M}} = \{\OLMathNumeral{1}, \OLMathNumeral{2}, \OLMathNumeral{3}\}$ و$\OLId{latin-ordinary}{s}(\OLId{latin-looped}{X}) = \{\OLMathNumeral{1}, \OLMathNumeral{2}\}$؛ أما إذا كان
  $\OLId{latin-ordinary}{s}(\OLId{latin-looped}{X}) = \{\OLMathNumeral{1}, \OLMathNumeral{2}, \OLMathNumeral{3}\} = \Domain{\OLId{latin-looped}{M}}$ فلا تُرضى.

  ولما امتنع إرضاء !!{formula} عند~$\OLId{latin-ordinary}{s}(\OLId{latin-looped}{X}) = \Domain{\OLId{latin-looped}{M}}$، كانت
  !!{sentence}
  \[
  \lforall[\OLId{latin-looped}{X}][\lexists[\OLId{latin-looped}{Y}][(\lexists[\OLId{latin-ordinary}{y}][\Atom{Y}{y}] \land
      \lforall[\OLId{latin-ordinary}{z}][(\Atom{X}{z} \liff \lnot \Atom{Y}{z})])]]
  \]
  غير مرضاة في كل حال. فمهما أخذنا !!{structure}~$\Struct{M}$،
  أتاح التعيين~$\OLId{latin-ordinary}{s}(\OLId{latin-looped}{X}) = \Domain{\OLId{latin-looped}{M}}$ إظهار كذب هذه !!{sentence}.
  وأما الجملة
  \[
  \lexists[\OLId{latin-looped}{X}][\lexists[\OLId{latin-looped}{Y}][(\lexists[\OLId{latin-ordinary}{y}][\Atom{Y}{y}] \land
      \lforall[\OLId{latin-ordinary}{z}][(\Atom{X}{z} \liff \lnot \Atom{Y}{z})])]]
  \]
  فإنها تُرضى مع أي تعيين~$\OLId{latin-ordinary}{s}$. وذلك أنه يوجد دائمًا
  $\OLId{latin-looped}{N} \subseteq \Domain{\OLId{latin-looped}{M}}$ مع~$\OLId{latin-looped}{N} \neq \Domain{\OLId{latin-looped}{M}}$؛ ومن ذلك
  $\OLId{latin-looped}{N} = \emptyset$.
\end{ex}

\begin{ex}
  معنى !!{sentence} الرتبة الثانية
  $\lforall[\OLId{latin-looped}{X}][\lforall[\OLId{latin-ordinary}{y}][\OLId{latin-looped}{X}(\OLId{latin-ordinary}{y})]]$ أن كل علاقة ذات~$\OLMathNumeral{1}$ موضع،
  أي كل خاصية، تصدق على كل فرد. وهذا ممتنع؛ ففي كل
  $\Struct{M}$ نستطيع أخذ تعيين متغيرات~$\OLId{latin-ordinary}{s}$ يحقق
  $\OLId{latin-ordinary}{s}(\OLId{latin-looped}{X}) = \emptyset$ و$\OLId{latin-ordinary}{s}(\OLId{latin-ordinary}{y}) = \OLId{latin-ordinary}{a} \in \Domain{\OLId{latin-looped}{M}}$، وحينئذ
  $\Sat/{\OLId{latin-looped}{M}}{\OLId{latin-looped}{X}(\OLId{latin-ordinary}{y})}[\OLId{latin-ordinary}{s}]$. ومن ثم تكون
  $!A \lif \lforall[\OLId{latin-looped}{X}][\lforall[\OLId{latin-ordinary}{y}][\OLId{latin-looped}{X}(\OLId{latin-ordinary}{y})]]$ مكافئة
  لـ$\lnot !A$ في منطق الرتبة الثانية؛ أي يكون
  $\Sat{\OLId{latin-looped}{M}}{!A \lif \lforall[\OLId{latin-looped}{X}][\lforall[\OLId{latin-ordinary}{y}][\OLId{latin-looped}{X}(\OLId{latin-ordinary}{y})]]}$ إذا وفقط إذا
  $\Sat{\OLId{latin-looped}{M}}{\lnot !A}$. فقد أمكن إذن تعريف~$\lnot$
  في هذا المنطق بوساطة~$\lforall$ و$\lif$.
\end{ex}

\begin{prob}
  بيّن أن بقية الروابط يمكن تعريفها أيضًا في منطق الرتبة الثانية
  بوساطة~$\lforall$ و$\lif$:
  \begin{enumerate}
    \item أثبت تكافؤ~$!A \land !B$ في هذا المنطق مع
    $\lforall[\OLId{latin-looped}{X}][((!A \lif (!B \lif \lforall[\OLId{latin-ordinary}{x}][\OLId{latin-looped}{X}(\OLId{latin-ordinary}{x})])) \lif
    \lforall[\OLId{latin-ordinary}{x}][\OLId{latin-looped}{X}(\OLId{latin-ordinary}{x})])]$.
    \item اطلب صيغة من الرتبة الثانية مقتصرة على~$\lforall$
    و$\lif$ تكافئ~$!A \lor !B$.
  \end{enumerate}
\end{prob}

\end{document}
